\chapter{End of Life Health}

\section*{Definition and Overview}
End of life health refers to the physical, emotional, and practical care of people in the period before death, and the support of their families. The goals of care change at the end of life, shifting from curing disease to relieving suffering, maintaining dignity, and honouring the person's wishes. Good end-of-life care addresses pain and symptoms, psychological and spiritual needs, and planning for the practicalities of death. This chapter explains what end-of-life care involves, how to plan for it, and the services that provide it.

\section*{Epidemiology}
Most Australians die of chronic, progressive conditions such as cancer, heart disease, dementia, and respiratory disease, often after a period of declining health. A growing proportion of people receive palliative care, but access varies, and many people die in hospital despite preferring home. Advance care planning rates remain low, meaning that many people do not have their wishes documented. End-of-life care is therefore an increasingly important part of health care as the population ages.

\section*{Aetiology and Pathophysiology}
The changes that occur at the end of life reflect the body's progressive shutdown.
\begin{itemize}
    \item \textbf{Underlying disease:} the illness itself, whether cancer, organ failure, or dementia, drives the final decline
    \item \textbf{Reduced energy:} the body conserves energy, with increasing sleep and reduced activity
    \item \textbf{Reduced intake:} appetite and thirst decline as the body needs less fuel
    \item \textbf{Circulatory changes:} blood is redirected to vital organs, making the hands and feet cool
    \item \textbf{Breathing changes:} breathing becomes irregular, with pauses and sometimes a rattling sound from secretions
    \item \textbf{Consciousness changes:} drowsiness increases and responsiveness fades, though hearing often remains
    \item \textbf{Pain and symptoms:} may be present and require active management
\end{itemize}

\section*{Clinical Features}
Families and carers should know what to expect in the final weeks, days, and hours.
\begin{itemize}
    \item Increasing drowsiness and withdrawal
    \item Loss of appetite and reduced drinking
    \item Difficulty swallowing medicines and fluids
    \item Confusion, restlessness, or agitation at times
    \item Irregular breathing and pauses
    \item Cool, mottled hands and feet
    \item Reduced urine output
    \item Changing levels of pain and comfort
    \item The person may hear even when unresponsive, so speaking gently is important
\end{itemize}

\section*{Differential Diagnosis}
Not every change at the end of life is expected, and some are treatable.
\begin{itemize}
    \item Pain, breathlessness, nausea, and agitation can all be controlled with appropriate care
    \item Infections, such as urinary or chest infections, may cause confusion and fever
    \item Medication side effects can cause drowsiness or restlessness
    \item Constipation and urinary retention cause distress and agitation
    \item Anxiety and spiritual distress are common and respond to support
\end{itemize}

\section*{When to Seek Medical Care}
Families should contact the palliative care team or GP for uncontrolled symptoms, concerns about medicines, or advice at any time. Good care plans include clear instructions on who to call and when.

\textbf{Call 000 (or local emergency services) for:}
\begin{itemize}
    \item Sudden severe breathlessness, chest pain, or collapse
    \item Sudden severe bleeding
    \item Seizures or fits
    \item Unexpected unconsciousness
    \item Any sudden change that needs immediate help
\end{itemize}

\section*{Diagnosis and Investigations}
Care at the end of life is assessed and planned rather than investigated.
\begin{itemize}
    \item \textbf{Advance care planning:} documenting values, preferences, and treatment wishes, including an advance care directive
    \item \textbf{Symptom assessment:} regular review of pain, breathing, nausea, and comfort
    \item \textbf{Medicine review:} stopping medicines that no longer help and continuing symptom relief
    \item \textbf{Discussion of resuscitation:} clarifying the person's wishes about CPR and life support
    \item \textbf{Minimal testing:} investigations are only used if they change the care plan
\end{itemize}

\section*{Management}
End-of-life care is delivered by a multidisciplinary team.
\begin{itemize}
    \item \textbf{Pain management:} regular and breakthrough medicines, including opioids, by the most appropriate route
    \item \textbf{Symptom control:} medicines for breathlessness, nausea, agitation, and secretions
    \item \textbf{Nursing care:} mouth care, positioning, and skin care for comfort
    \item \textbf{Emotional and spiritual support:} counselling, chaplaincy, and time for family
    \item \textbf{Family support:} practical help, information, and bereavement care
    \item \textbf{Place of care:} home, hospice, or hospital, according to wishes and needs
    \item \textbf{After death:} clear guidance on what happens next
\end{itemize}

\section*{Complications}
\begin{itemize}
    \item Uncontrolled pain and symptoms if care is inadequate
    \item Distress and complicated grief in families
    \item Unmet wishes about place or manner of death when planning is absent
    \item Carer exhaustion and burnout
    \item Financial and practical strain on families
\end{itemize}

\section*{Prognosis and Outcomes}
With good palliative care, most people die peacefully and free of pain. Predicting the exact time of death is difficult, but the final phase follows a recognisable pattern that families can be prepared for. Families who are involved, informed, and supported report better experiences, and knowing that the person's wishes were honoured eases grief.

\section*{Prevention and Self-Care}
\begin{itemize}
    \item Document your wishes in an advance care directive and review it regularly
    \item Appoint a substitute decision-maker and tell them what matters to you
    \item Talk to family about preferences for care and place of death
    \item Learn about palliative care services before they are needed
    \item Keep important documents accessible: will, power of attorney, and directive
    \item Ask the care team anything; there are no silly questions
    \item Support family members to rest and care for themselves
\end{itemize}

\section*{Sources and Further Reading}
The following reputable sources provide further information for healthcare students and patients seeking reliable, current advice about end-of-life health.
\begin{itemize}
    \item Palliative Care Australia. \textit{End-of-life care information and support.} https://palliativecare.org.au/
    \item Advance Care Planning Australia. \textit{Planning for end of life.} https://www.advancecareplanning.org.au/
    \item CareSearch. \textit{Palliative care and end-of-life resources.} https://www.caresearch.com.au/
    \item NHS. \textit{End of life care.} https://www.nhs.uk/conditions/end-of-life-care/
    \item Dying to Talk. \textit{Conversation resources.} https://www.dyingtotalk.org.au/
    \item MedlinePlus. \textit{End of life issues.} https://medlineplus.gov/endoflifeissues.html
\end{itemize}
